\section{Methods}\label{sec:model}

\subsection{A two-land model with conserved area and a separate quality state}\label{sec:equations}

We divide total bioproductive area $A_{tot}$ among three land classes. Fast-provisioning land $A_f$ represents cropland and intensively managed systems. Ecological-capital land $A_c$ represents forest, soil, wetland, and long-rotation systems as one effective class, rather than as a single real land type. Reserve land $A_r$ is fallow, retired, or not yet allocated. The areas always sum to $A_{tot}$. Capital land also has an uncalibrated, model-defined quality index $q\in K_q=[0,q_{\max}]$ that affects its annual account contribution; $q$ is not an area or an observed ecological variable. We impose $q_{\max}=1$ as an admissibility bound, not as an estimated ecological threshold. The raw quality-regeneration rate is logistic:
\begin{equation}\label{eq:Gc}
G_c(q)=\rho_c\,q\left(1-\frac{q}{q_{\max}}\right),
\end{equation}
We calculate the annual value of capital land per hectare as its direct yield plus a term that depends on quality growth:
\begin{equation}\label{eq:Ups}
\Ups(q)=b_c+b_{G,c}\,G_c(q),
\end{equation}
Here $b_c$ is the direct yield, and $b_{G,c}$ converts quality growth into added annual value. Total biocapacity is the sum of the fast-land and capital-land contributions
\begin{equation}\label{eq:B}
B=Y_f+Y_c=b_f A_f+\Ups(q)\,A_c ,
\end{equation}
The resulting account amount $B$ is in gha, with the annual accounting basis implicit \citep{GFNStandards2009,GFNGuidebook2019}. Demand is $E=eP$: population $P$ multiplied by the constant per-person footprint $e$ (gha per person; $\mathrm{cap}$ denotes one person in the equations and tables). The shortfall $S$ is the part of demand not met by the shares of the two land contributions available to people:
\begin{equation}\label{eq:S}
S=\pos{E-\sigma_f Y_f-\sigma_c Y_c},
\end{equation}
Here $\pos{x}=\max(x,0)$, and $\sigma_f=\sigma_c=1$ in the baseline. When a shortfall remains, the model converts capital land into fast-provisioning land at a finite rate:
\begin{equation}\label{eq:uc}
u_c=\frac{1}{\tconv}\min\!\left(\frac{S}{b_f},\;A_c-\Acmin\right)\qquad[\mathrm{ha\,yr^{-1}}],
\end{equation}
Here $S/b_f$ is the area needed to cover the shortfall at the fast-land yield, $\Acmin$ is the minimum capital-land area allowed, and $\tconv$ sets how quickly conversion occurs. Land moves between classes through one-to-one transfers:
\begin{align}
\dot A_c&=-u_c+R_{rc},\label{eq:Ac}\\
\dot A_f&=u_c+\mu A_r-\eta_f A_f,\label{eq:Af}\\
\dot A_r&=-\mu A_r+\eta_f A_f-R_{rc},\label{eq:Ar}
\end{align}
Here $\mu A_r$ moves reserve land into fast use (for example, through replanting), and $\eta_f A_f$ moves fast land back into reserve. The only route from reserve back to ecological-capital land is the surplus-gated restoration flow $R_{rc}$:
\begin{equation}\label{eq:Rrc}
R_{rc}=\chi_r\,A_r\,\pos{B-E},
\end{equation}
These transfer equations conserve total area, so reserve land is the remainder after $A_f$ and $A_c$ are set. The delayed quality law projects its regeneration vector field onto the tangent cone of this admissible interval. The model also tracks accumulated shortfall as debt, which lowers fast-land yield. Population follows a logistic response to delayed carrying capacity $K(t-\tau_p)$, with $K=\max(B/e,K_{\min})$; the delay is in the capacity input, not in the population state in its crowding term:
\begin{align}
\dot q&=\Pi_{T_{K_q}(q(t))}G_c\big(q(t-\tau_g)\big),\label{eq:q}\\
\dot P&=rP\left(1-\frac{P}{K(t-\tau_p)}\right),\qquad K=\max(B/e,\,K_{\min}),\label{eq:P}\\
b_f&=\big(b_{f0}+T_b(t)\big)\,e^{-\alpha D},\label{eq:bf}\\
\dot D&=\pos{E-B}-\eta D,\label{eq:D}\\
T_b(t)&=\omega_{Tb}\,\frac{\Delta b}{1+e^{-\kappa_w(t-t_{wave})}},\qquad \omega_{Tb}\in\{0,1\}.\label{eq:Tb}
\end{align}
Here $T_{K_q}(q)$ is the tangent cone to $K_q$ at the current state, and $\Pi$ denotes Euclidean projection onto that cone. We interpret $q$ as an absolutely continuous solution satisfying the projected equation almost everywhere, with continuous initial history $q(s)=\phi(s)\in K_q$ on $[-\tau_g,0]$. Since $G_c(x)\ge0$ for $x\in K_q$, the projected field agrees with the raw delayed regeneration rate below the ceiling and sets the outward velocity to zero at $q=q_{\max}$; at $q=0$ it has no negative outward component.

The model keeps area, annual account amounts, and accumulated shortfall distinct. Land areas $A_f$, $A_c$, and $A_r$ are in hectares; $B$, $E$, and $S$ are annual account amounts in gha, with the annual basis implicit. $D$ accumulates the shortfall through time, so it is measured in gha\,yr, while $\dot D$ and the repayment term $\eta D$ are in gha. This debt is an internal model quantity, not an area or an item reported by the NFA.

The optional technology boost $T_b$ raises fast-land yield gradually, following a bounded logistic curve. It is off in the baseline, main figures, and tables; it is switched on only in the matched comparison in SI~\S S8. The model tracks fast area $A_f$, capital area $A_c$, capital quality $q$, population $P$, and accumulated shortfall $D$. Timber liquidation and direct restoration from fast land to capital land are possible extensions, but are off in all reported runs (SI~\S S1).

The conversion switch is exact: it turns off when the shortfall reaches zero, and land moves at a finite rate set by $\tconv$. Changes in $q$ alter value per hectare without moving land. Reserve replanting can maintain fast land when conversion is inactive; without replanting, retirement would drive fast-land area to zero at a no-conversion equilibrium.

The model has two reinforcing feedback loops. First, if a hectare of fast land contributes more annual biocapacity than a hectare of capital land ($b_f>\Ups(q)$), converting land raises $B$ at that moment. More biocapacity can support a larger population; after the demographic delay, higher population can raise demand and create another shortfall. Other land flows and quality changes also affect $B$, so this conversion effect alone does not determine whether total biocapacity rises or falls.

Second, a shortfall adds to debt $D$, while repayment reduces it. More debt lowers fast-land yield $b_f$, which can reduce $B$ and deepen the shortfall. With fixed demand and a continuing shortfall, debt approaches a finite level.

\subsection{Parameters and their status}\label{sec:params}

Table~\ref{tab:params} lists the baseline values and their roles. The model groups forests, soils, wetlands, and other slow-renewing assets into one effective capital-land class, although they differ in area, condition, yield, and recovery time. The quality index $q$ is uncalibrated, $q_{\max}=1$ is imposed, and the baseline values of $\rho_c$ and $\tau_g$ are scenario settings rather than estimates. Published records show heterogeneous recovery times: forest biomass reaches about half of its reference in a median of $30$\,yr, collapsed fish stocks in about $5.6$\,yr, and soil organic carbon returns toward a new balance over roughly $17$--$23$\,yr (SI~\S S5.5) \citep{Poorter2016,HutchingsReynolds2004,Neubauer2013,Poeplau2011}. These records provide context for future ecosystem-specific measurement and modeling only; they do not validate the present quality equation or identify $q$, $\rho_c$, or $\tau_g$.

The NFA cropland account value divided by cropland area, and the separate cereal-yield/cropland-area proxy, describe account data. Neither is a fitted fast-land yield $b_f$ or a direct measure of ecological capital (SI~\S S5). Their boundaries do not establish a one-to-one match with model fast land. The baseline has $b_f-\Ups(q)>0$, so a transfer raises biocapacity at that moment; this is a conditional model setting, not an empirical prediction.

\begin{table}[!htbp]
\centering\footnotesize
\renewcommand{\arraystretch}{0.95}
\caption{Baseline parameter values and evidence status. Representative values are effective, lumped quantities; the NFA account/area quotient is descriptive, not a physical-yield measurement. Extension parameters are active only in scenarios specified in this table or the SI. Symbols and units are listed in SI~\S S1.}
\label{tab:params}
\setlength{\tabcolsep}{4pt}
\begin{tabularx}{\textwidth}{@{}>{\raggedright\arraybackslash}p{0.13\textwidth} >{\raggedright\arraybackslash}p{0.2\textwidth} >{\raggedright\arraybackslash}X >{\raggedright\arraybackslash}p{0.26\textwidth}@{}}
\toprule
Parameter & Value & Meaning & Status \\
\midrule
$\rho_c$ & $0.08$\,yr$^{-1}$ & quality-regeneration scenario rate & scenario setting; not estimated \\
$b_{f0}$ & $0.85$\,gha\,ha$^{-1}$ & base fast-land yield per hectare & representative (2026 NFA account value/area quotient $\approx1.045$ in 1961; not a physical-yield measure) \\
$b_c$ & $0.05$\,gha\,ha$^{-1}$ & capital-land direct yield & representative \\
$b_{G,c}$ & $0.8$\,gha\,yr\,ha$^{-1}$ & regeneration-value coefficient for capital land & representative \\
$q_{\max}$ & $1.0$ (dimensionless) & imposed quality admissibility bound & hard state constraint; not estimated \\
$\mu$ & $0.06$\,yr$^{-1}$ & cropland maintenance (reserve$\to$fast) & representative \\
$\eta_f$ & $0.05$\,yr$^{-1}$ & fast-land retirement (fast$\to$reserve) & representative \\
$\tconv$ & $1.0$\,yr & conversion time constant & representative \\
$\alpha$ & $0.03$\,gha$^{-1}$\,yr$^{-1}$ & degradation of $b_f$ per unit debt & representative \\
$\eta$ & $0.05$\,yr$^{-1}$ & debt-repayment rate ($>0$) & representative \\
$e$ & $0.55$\,gha\,cap$^{-1}$ & per-capita footprint (constant) & representative \\
$r$ & $0.02$\,yr$^{-1}$ & intrinsic population growth & representative \\
$\tau_g,\ \tau_p$ & $20,\ 25$\,yr & regeneration and demographic delays & scenario settings; $\tau_p$ lumps institutional response \\
$A_{tot}$ & $2.5$\,ha & total bioproductive area & normalization \\
$\sigma_f,\ \sigma_c$ & $1,\ 1$ (dimensionless) & human-available shares of annual account amounts & baseline ($<1$: reservation extension) \\
$\chi_r$ & $0.10$\,gha$^{-1}$\,yr$^{-1}$ & surplus-gated restoration coefficient & on with endogenous demand; off under frozen demand \\
$\Acmin$, $A_r^{\min}$,\newline $K_{\min}$ & $0.05$\,ha, $0.10$\,ha,\newline $0.02$\,cap & typed floors & domain constraints \\
$\Delta b$, $\kappa_w$,\newline $t_{wave}$ & $1.0$\,gha\,ha$^{-1}$, $0.1$\,yr$^{-1}$,\newline $40$\,yr & technology wave & off in baseline; active only in the matched SI S8 experiment \\
\bottomrule
\end{tabularx}
\end{table}

\subsection{Data and empirical measures}\label{sec:data}

We use the official National Ecological Footprint and Biocapacity Accounts, Edition 2026, results and metadata version 1.0 \citep{Lo2026NFA}. The main analysis covers the complete world series from 1961 to 2023; estimated values for 2024--2025 are excluded. The 1961--2022 sensitivity uses the same edition and population series, so both windows share one data vintage.

From the workbook's \texttt{world\_data} sheet, we use \texttt{biocap\_total\_gha} for total biocapacity $B$, \texttt{efp\_total\_gha} for the Global Ecological Footprint of Production $E_{\mathrm{prod}}$, and \texttt{pop\_1000\_persons} for UN PopStat population. We convert population from thousands to persons and calculate per-person values from the unrounded totals, rather than from the workbook's rounded per-person columns. Ratios, indices, and logarithmic changes have no units.

The component series retain the workbook labels Cropland, Grazing land, Forest biocapacity, Fishing grounds, and Built-up land. These are NFA account values, not measurements of biomass, soil condition, wetland extent, fish biomass, or the model's effective area $A_c$. Built-up land is an NFA accounting class and is not ecological capital. Accordingly, the land-type movements are used as accounting-composition signals only. Variable definitions and endpoint calculations are summarized in SI~\S S5.

As a separate check, we combine OWID's FAOSTAT-derived world cropland-area series (arable land plus permanent crops) with cereal yield \citep{OWIDCereal2026,OWIDCropland2026}. Cereal yield is a limited indicator of crop productivity, not a measure of all-crop output or an NFA yield factor. We scale the proxy to match the NFA cropland share of total biocapacity in 1961, then calculate $X_t=\hat b_{f,t}A_{f,t}$ and the remainder $C_t=B_t-X_t$. This is an assumed account proxy, not a measured split of physical capital. The remainder can include differences in boundaries, categories, and calibration; it is not $\Upsilon(q)A_c$. The proxy and its calibration sensitivity are reported in SI~\S S5. We infer no physical-stock trend from it.

The aggregate-versus-per-capita comparison uses the exact identity $\Delta\ln B=\Delta\ln(B/P)+\Delta\ln P$. This extensive/intensive split is an accounting identity, not a causal decomposition of biocapacity production.

\subsection{Numerical solution}\label{sec:numerics}

We use the method of steps and an explicit first-order Euler scheme. For ordinary runs, pre-run histories of $q$ and $B$ are constant at their initial values over the longest active delay; delayed values are evaluated by linear interpolation. In the wave experiment, both runs share the no-wave pre-run $B$ history, while the active scenario changes $B$ at $t=0$. At each step, the solver reads the state at the previous grid point, evaluates the time-dependent technology term and delayed inputs at the next grid time, computes the flows and derivatives, projects the quality velocity onto the tangent cone before the Euler advance, then applies the interval clip as a finite-step safeguard, clips the other bounded states, computes reserve area as the area-conservation residual, and stores the updated $B$ history. The demonstration run and drawdown-depth sweeps use $\Delta t=0.05$ and $0.1$\,yr, respectively; the wave runs use $0.05$ and selected $0.025$\,yr steps. The positive-part threshold is exact. The projected specification and the finite-step update were checked separately: at the same $\Delta t=0.05$\,yr grid over $T=100$\,yr with the Figure~1 baseline initial state, the archived v45 candidate's post-Euler clip and the explicit tangent-cone Euler update produced identical state arrays for $\tau_g=5,20,40$\,yr. For the baseline quality history, the continuous projected solution at $\tau_g=20$ or $40$ is $q(t)=\min(1,0.9+0.0072t)$ and reaches the ceiling at $125/9\approx13.8889$\,yr; the piecewise-linear discrete path's maximum error fell from $7.1\times10^{-5}$ at $\Delta t=0.1$\,yr to $8.8\times10^{-6}$ at $0.0125$\,yr. In the full frozen-demand Figure~1 run, halving the step from $0.05$ to $0.025$\,yr reduced maximum absolute differences versus the $0.003125$\,yr reference from $0.00221$ to $0.00102$\,ha for $A_c$, $0.00180$ to $0.000831$\,gha for $B$, $0.000252$ to $0.000118$\,cap for $P$, and $0.00200$ to $0.000924$\,gha\,yr for $D$. These are checks of the local v45-derived implementation, not proof that every archived C10 output was generated by this exact run path; the historical provenance is assessed separately. Floor-proximity values remain tolerance- and resolution-dependent, so they are reported as bands rather than sharp thresholds.
